\documentclass[uplatex,dvipdfmx,a4paper,11pt]{jsarticle}
\input{../../../templates/kaisetsu.tex}
\kaisetsuprint{量子力学1　第10章「自由粒子の波束（群速度・波束の広がり）」}{2026-11-21}
%% kaisetsu-moto	problems\量子力学1\shou10\量子力学1-10.tex	kiroku\kaisetsu\量子力学1\量子力学1-10_2026-11-21_komento.tex
\begin{document}
\matome
\begin{saitenichiran}
  \hyougyou{\hyperlink{bt-m-075-dispersion}{1}}{L1・計算}{◎◎○}{A}{(3)記述}
\end{saitenichiran}
\begin{daiji}
  \item 分散関係の定数項（静止エネルギー $mc^{2}/\hbar$）は波動関数全体に共通の位相 $\ee^{-\ii mc^{2}t/\hbar}$ を掛けるだけで $\abs{\psi}^{2}$ を変えないので，観測される粒子の運動（包絡の速さ $v_{\mathrm{g}}$）には効かない。$v_{\mathrm{p}}$ はエネルギーの原点の取り方で変わる量で，粒子の速さではない（問題1 (3)）。
\end{daiji}
\begin{maemo}
  \item 位相速度の扱い：第2章 問題5（2026-10-29）で位相速度を $\hbar k/m$ と取り違えた。位相速度は $\omega/k$ で，エネルギーの原点をずらすと変わる。物理的な速さとして使えるのは $\dd\omega/\dd k$ の方である，と毎回確かめる。
\end{maemo}

\mondaihen
\level{1}
\mondaibangou{1}
\begin{mondai}{075-dispersion}{計算}{相対論的な物質波の分散関係での位相速度・群速度と，群速度が粒子の速さを表すこと（エネルギーの原点への依存）}
\begin{teisuu}
  換算プランク定数 & $\hbar$ & \qty{1.055e-34}{\joule\second}\\
  電子の質量 & $\melec$ & \qty{9.109e-31}{\kilogram}\\
\end{teisuu}
相対論的な自由粒子のエネルギー $E=\sqrt{p^{2}c^{2}+m^{2}c^{4}}$（第1章）に $E=\hbar\omega$，$p=\hbar k$ を入れると，物質波の分散関係
  \[ \omega(k)=c\sqrt{k^{2}+\mu^{2}},\qquad \mu=\frac{mc}{\hbar} \]
  を得る（$k>0$）。
  \begin{shoumon}
    \item 位相速度 $v_{\mathrm{p}}=\omega/k$ と群速度 $v_{\mathrm{g}}=\dd\omega/\dd k$ を求めよ。積 $v_{\mathrm{p}}v_{\mathrm{g}}$ の値と，$v_{\mathrm{p}}$，$v_{\mathrm{g}}$ のそれぞれが $c$ より大きいか小さいかを答えよ。
    \item 速さ $v$ の相対論的な粒子では $p=Ev/c^{2}$ である。$v_{\mathrm{g}}$ が粒子の速さ $v$ に一致することを示せ。
    \item $k\ll\mu$ のとき，$\omega(k)$ を $k^{2}$ の1次まで展開し，$v_{\mathrm{g}}$ と $v_{\mathrm{p}}$ の近似式を求めよ。非相対論の分散関係 $\omega=\hbar k^{2}/(2m)$（第2章）から求めた $v_{\mathrm{g}}$，$v_{\mathrm{p}}$ と比べ，一致するのはどちらかと，その理由を述べよ。
  \end{shoumon}

\begin{kotaeran}
  \machigai{3}{○}{記述} 展開 $\omega\approx c\mu+ck^{2}/(2\mu)=mc^{2}/\hbar+\hbar k^{2}/(2m)$，$v_{\mathrm{g}}\approx\hbar k/m$，$v_{\mathrm{p}}\approx mc^{2}/(\hbar k)+\hbar k/(2m)$ と，非相対論の値と比べて $v_{\mathrm{g}}$ が一致することは正しい。理由が「$v_{\mathrm{p}}$ は静止エネルギー分ずれる」だけで，問われている「一致する理由」の中心が書かれていない。(i) 定数項 $mc^{2}/\hbar$ は $k$ で微分すると消えるので $v_{\mathrm{g}}$ には効かず，$\omega/k$ には $mc^{2}/(\hbar k)$ として残る。(ii) 定数項は全体に共通の位相を掛けるだけで $\abs{\psi}^{2}$ を変えないので，観測される運動は変わらない。この2点まで書く。
\end{kotaeran}
\begin{rikai}
  \item $\omega\to\omega+\omega_{0}$ と定数をずらすと $\psi\to\psi\,\ee^{-\ii\omega_{0}t}$ となり，$\abs{\psi}^{2}$ は変わらない。観測にかかるのは $\omega$ の $k$ による変化だけで，$\dd\omega/\dd k$ はそれだけで決まる。
  \item 「どちらが一致するか，その理由」の問いには，式の上の理由（微分で定数が消える）と物理的な理由（共通の位相は観測されない）の両方を書く。
\end{rikai}
\end{mondai}

\kaitouhen
\level{1}
\kaisetsubangou{1}
\begin{kaisetsu}{075-dispersion}
\houshin 分散関係を微分して群速度を求め，$p=\hbar k$，$E=\hbar\omega$ で粒子の速さと比べる。(3) は展開で現れる定数項が $\dd\omega/\dd k$ には効かず，$\omega/k$ には効くことに注目する。
\keisan \shoumonkai{1} 定義より
\[ v_{\mathrm{p}}=\frac{\omega}{k}=\frac{c\sqrt{k^{2}+\mu^{2}}}{k}=c\sqrt{1+\frac{\mu^{2}}{k^{2}}}>c \]
である。合成関数の微分で
\[ v_{\mathrm{g}}=\dv{\omega}{k}=c\cdot\frac{2k}{2\sqrt{k^{2}+\mu^{2}}}=\frac{ck}{\sqrt{k^{2}+\mu^{2}}}<c \]
となる（$\mu>0$ なので分母は $k$ より大きい）。積は
\[ v_{\mathrm{p}}v_{\mathrm{g}}=\frac{c\sqrt{k^{2}+\mu^{2}}}{k}\cdot\frac{ck}{\sqrt{k^{2}+\mu^{2}}}=c^{2} \]
である。別解として，$\omega^{2}=c^{2}(k^{2}+\mu^{2})$ を $k$ で微分すると $2\omega\,\dd\omega/\dd k=2c^{2}k$ で，両辺を $2k$ で割ると $(\omega/k)(\dd\omega/\dd k)=c^{2}$ となる。

\shoumonkai{2} $p=Ev/c^{2}$ より $v=pc^{2}/E$ である。$p=\hbar k$，$E=\hbar\omega$ を入れると
\[ v=\frac{\hbar kc^{2}}{\hbar\omega}=\frac{kc^{2}}{c\sqrt{k^{2}+\mu^{2}}}=\frac{ck}{\sqrt{k^{2}+\mu^{2}}}=v_{\mathrm{g}} \]
となる。

\shoumonkai{3} $k\ll\mu$ では $\sqrt{k^{2}+\mu^{2}}=\mu\sqrt{1+k^{2}/\mu^{2}}\approx\mu(1+k^{2}/(2\mu^{2}))$ なので，$\mu=mc/\hbar$ を戻すと
\[ \omega(k)\approx c\mu+\frac{ck^{2}}{2\mu}=\frac{mc^{2}}{\hbar}+\frac{\hbar k^{2}}{2m} \]
である。これから
\[ v_{\mathrm{g}}\approx\frac{\hbar k}{m},\qquad v_{\mathrm{p}}\approx\frac{mc^{2}}{\hbar k}+\frac{\hbar k}{2m} \]
となる。非相対論の $\omega=\hbar k^{2}/(2m)$ からは $v_{\mathrm{g}}=\hbar k/m$，$v_{\mathrm{p}}=\hbar k/(2m)$ なので，$v_{\mathrm{g}}$ は一致するが $v_{\mathrm{p}}$ は一致しない。違いは定数項 $\omega_{0}=mc^{2}/\hbar$（静止エネルギー，すなわちエネルギーの原点の取り方）だけから来る。定数は $k$ で微分すると消えるので $v_{\mathrm{g}}$ には効かないが，$\omega/k$ には $\omega_{0}/k$ として残る。定数項は波動関数全体に共通の位相 $\ee^{-\ii\omega_{0}t}$ をかけるだけで $\abs{\psi}^{2}$ を変えないので，観測される粒子の運動（包絡の速さ $v_{\mathrm{g}}$）は変わらない。$v_{\mathrm{p}}$ はエネルギーの原点に依存し，粒子の速さを表す量ではない。
\kotae{\shoumonkai{1} $v_{\mathrm{p}}=c\sqrt{k^{2}+\mu^{2}}/k>c$，$v_{\mathrm{g}}=ck/\sqrt{k^{2}+\mu^{2}}<c$，$v_{\mathrm{p}}v_{\mathrm{g}}=c^{2}$ \shoumonkai{2} $v=pc^{2}/E=ck/\sqrt{k^{2}+\mu^{2}}=v_{\mathrm{g}}$ \shoumonkai{3} $v_{\mathrm{g}}\approx\hbar k/m$（一致する），$v_{\mathrm{p}}\approx mc^{2}/(\hbar k)+\hbar k/(2m)$（静止エネルギーの分だけ一致しない）}
\begin{ronten}
  \item \shoumonkai{1} $\mu>0$ なので $\sqrt{k^{2}+\mu^{2}}>k$ であることを，$v_{\mathrm{p}}>c$，$v_{\mathrm{g}}<c$ の根拠として書く。
  \item \shoumonkai{3} $k\ll\mu$ なので $k^{2}/\mu^{2}$ を小さい量として平方根を展開したことを書く。
  \item \shoumonkai{3} 展開で現れる定数項が $k$ で微分すると消えるので $v_{\mathrm{g}}$ には効かず，$\omega/k$ には残ることを述べる。
  \item \shoumonkai{3} 定数項は全体に共通の位相をかけるだけで $\abs{\psi}^{2}$ を変えないので，観測される粒子の運動は変わらないことを書く。
\end{ronten}
\kenzan $\mu\to0$（光子）では $v_{\mathrm{p}}=v_{\mathrm{g}}=c$，$k\to\infty$（超相対論）でも両方とも $c$ に近づく。$c\mu$ は $\unit{\metre\per\second}\cdot\unit{\per\metre}=\unit{\per\second}$ で角振動数の次元をもつ。(3) の $v_{\mathrm{p}}$ の主要項 $c\mu/k$ は $k\ll\mu$ で $c$ よりはるかに大きく，(1) の $v_{\mathrm{p}}>c$ と合う。
\tsumazuki $v_{\mathrm{p}}>c$ は相対論に反しない。位相が一定の点は情報やエネルギーを運ばないからである。
\end{kaisetsu}

\end{document}
